% =====================================================================
% CSE 341 -- Concepts of Programming Languages
% Project D3 (Part 1) -- Example Programs & Test Report for ANS
% Students: 220104004926 & 220104004931
% Compile with: tectonic 220104004926-220104004931_D3_P1.tex
% =====================================================================
\documentclass[10pt,a4paper]{article}

\usepackage[utf8]{inputenc}
\usepackage[T1]{fontenc}
\usepackage[english]{babel}
\usepackage[a4paper,margin=1.65cm]{geometry}
\usepackage{microtype}
\usepackage{xcolor}
\usepackage{enumitem}
\usepackage{upquote}
\usepackage{listings}
\usepackage[hidelinks]{hyperref}

\definecolor{ansblue}{HTML}{1E3A8A}
\definecolor{codebg}{HTML}{F8FAFC}
\definecolor{codeborder}{HTML}{CBD5E1}

\hypersetup{
  pdftitle={D3 - Example Programs and Test Report},
  pdfauthor={Yusuf Hamza Çelebi and Enes Özkan}
}

\setlength{\parindent}{0pt}
\setlength{\parskip}{4pt}
\setlist[itemize]{leftmargin=*, itemsep=2pt, topsep=2pt}
\setlist[description]{leftmargin=1.9cm, style=nextline, itemsep=3pt, topsep=2pt}

\lstdefinestyle{terminal}{
  basicstyle=\ttfamily\scriptsize,
  backgroundcolor=\color{codebg},
  frame=single,
  rulecolor=\color{codeborder},
  framerule=0.3pt,
  xleftmargin=3pt,
  xrightmargin=3pt,
  aboveskip=5pt,
  belowskip=5pt,
  columns=fullflexible,
  keepspaces=true,
  upquote=true,
  showstringspaces=false,
  breaklines=true,
  breakatwhitespace=false
}

\newcommand{\code}[1]{\texttt{#1}}

\begin{document}

\begin{center}
  {\Large\bfseries D3 -- Example Programs \& Test Report}\\[2pt]
  {\large ANS -- Atomic Notation Script}\\[4pt]
  Yusuf Hamza Çelebi (220104004926) and Enes Özkan (220104004931)\\
  CSE 341 -- Programming Languages, Spring 2026
\end{center}

\section*{Purpose and Parser Under Test}

This report documents the Part 1 D3 tests for the ANS lexer and parser. The parser under test is the C implementation in \code{D2/}, run as \code{D2/ans <file>}. A valid program must parse with exit code \code{0}; a malformed program must be rejected with a useful diagnostic containing the file name, line number, column, and message.

\section*{Three Valid Programs}

The full valid source files are submitted in \code{D3/programs/}. The parser accepts all three programs silently, which is the expected success behavior for \code{D2/ans <file>}.

\begin{itemize}
  \item \textbf{\code{molar\_mass\_demo.ans} (40 lines).}
  Exercises \code{element} and \code{compound} declarations, functions, typed parameters, list literals, \code{for}, \code{if}/\code{elif}/\code{else}, member access, index access, and function calls. Result: accepted, exit \code{0}.

  \item \textbf{\code{combustion\_analysis.ans} (52 lines).}
  Exercises a \code{reaction} declaration, formula literals, \code{check = balanced}, exothermic energy quantity, typed \code{reaction}/\code{compound} parameters, nested calls, and augmented assignment. Result: accepted, exit \code{0}.

  \item \textbf{\code{multi\_reaction\_survey.ans} (70 lines).}
  Exercises multiple reactions, formula-literal and identifier species, compound lists, loop counters, relational expressions, character and string literals, and repeated function calls. Result: accepted, exit \code{0}.
\end{itemize}

Verification run:
\begin{lstlisting}[style=terminal]
D2/ans D3/programs/combustion_analysis.ans     OK
D2/ans D3/programs/molar_mass_demo.ans         OK
D2/ans D3/programs/multi_reaction_survey.ans   OK
\end{lstlisting}

Together, these examples cover the required DSL surface: primitive literals (\code{number}, \code{string}, \code{char}, contextual words), structured chemistry declarations (\code{element}, \code{compound}, \code{reaction}), functions with parameters, expression precedence, control flow, list traversal, and the chemistry-specific reaction syntax.

\section*{Five Malformed Programs}

The full malformed source files are submitted in \code{D3/malformed/}. Each malformed file is intentionally small so the reported line number points directly to the syntax or lexical problem.

\begin{itemize}
  \item \textbf{\code{err1\_missing\_semicolon.ans}.}
  Missing \code{;} after an \code{element} declaration. The parser reports the next token, \code{let}, at line 7 where the semicolon was required.

  \item \textbf{\code{err2\_unknown\_field.ans}.}
  Invalid field name \code{protons} inside an \code{element} body. The parser reports the unknown field at line 5, then a recovery error at the following \code{=}.

  \item \textbf{\code{err3\_bad\_identifier.ans}.}
  Illegal character \code{@} in an expression. The lexer reports the illegal character with line and column, then the parser reports the missing semicolon.

  \item \textbf{\code{err4\_unmatched\_brace.ans}.}
  Missing closing brace for a \code{for} block inside a function. The parser reports a syntax error near the premature \code{;}, followed by EOF recovery diagnostics.

  \item \textbf{\code{err5\_unterminated\_string.ans}.}
  String literal reaches a newline before a closing quote. The lexer reports the unterminated string at line 5, then parser recovery reports the broken compound body.
\end{itemize}

Actual parser output:
\begin{lstlisting}[style=terminal]
=== D3/malformed/err1_missing_semicolon.ans ===
D3/malformed/err1_missing_semicolon.ans:7:1: error: expected ';' but got let

=== D3/malformed/err2_unknown_field.ans ===
D3/malformed/err2_unknown_field.ans:5:5: error: unknown element field 'protons'
D3/malformed/err2_unknown_field.ans:5:13: error: expected '}' but got =

=== D3/malformed/err3_bad_identifier.ans ===
D3/malformed/err3_bad_identifier.ans:3:18: error: unexpected character '@'
D3/malformed/err3_bad_identifier.ans:3:18: error: expected ';' but got @

=== D3/malformed/err4_unmatched_brace.ans ===
D3/malformed/err4_unmatched_brace.ans:8:2: error: unexpected token ; in expression
D3/malformed/err4_unmatched_brace.ans:9:1: error: expected ';' but got
D3/malformed/err4_unmatched_brace.ans:9:1: error: expected '}' but got
D3/malformed/err4_unmatched_brace.ans:9:1: error: expected ';' but got

=== D3/malformed/err5_unterminated_string.ans ===
D3/malformed/err5_unterminated_string.ans:5:13: error: unterminated string literal
D3/malformed/err5_unterminated_string.ans:5:13: error: expected STRING but got "
D3/malformed/err5_unterminated_string.ans:6:2: error: expected '}' but got ;
D3/malformed/err5_unterminated_string.ans:7:1: error: expected ';' but got
\end{lstlisting}

\section*{Conclusion}

The valid examples demonstrate that the Part 1 parser accepts non-trivial ANS programs using declarations, expressions, control structures, functions, and the chemistry-specific reaction construct. The malformed examples show that lexical errors, missing separators, invalid declaration fields, unmatched braces, and unterminated strings are rejected with line-numbered, actionable diagnostics.

\end{document}
